\documentclass[10pt]{article}

\usepackage[a4paper,margin=2cm]{geometry}
\usepackage[utf8]{inputenc}
\usepackage{amsmath}
\usepackage{amssymb}

\usepackage{longtable}
\usepackage{array}
\usepackage{hyperref}
\usepackage{xcolor}
\usepackage{booktabs}

\hypersetup{
    colorlinks=true,
    urlcolor=blue,
    linkcolor=black
}

\newcolumntype{P}[1]{>{\raggedright\arraybackslash}p{#1}}

\title{Glossary of Abbreviations for S+SSPR 2026}
\author{A. Schlapbach}
\date{}

\begin{document}


\begin{longtable}{P{2cm} P{3cm} P{9cm} P{2cm}}
\toprule
\textbf{Abbrev.} & \textbf{Full Term} & \textbf{Description} & \textbf{Reference} \\
\midrule
\endhead

PCA & Principal Component Analysis &
A linear dimensionality-reduction technique that projects correlated variables onto orthogonal axes (principal components) ordered by explained variance. It works by eigendecomposing the covariance matrix of the standardized data; here it compresses 7 raw airline variables into 3--4 components while retaining 87.8--95.6\% of total variance. &
\href{https://en.wikipedia.org/wiki/Principal_component_analysis}{Wikipedia} \\

pc & Principal Component &
A single output axis of a PCA decomposition (pc1, pc2, pc3, pc4), ranked by the share of variance it explains. Each is a linear combination of the original 7 variables, with ``loadings'' indicating how strongly each original variable contributes; the paper interprets pc1--pc4 substantively (growth, profit--wage tension, fuel shocks, COVID yield collapse). &
\href{https://en.wikipedia.org/wiki/Principal_component_analysis#Details}{Wikipedia} \\

LF & Load Factor &
An airline-industry KPI equal to revenue passenger miles divided by available seat miles, expressed as a percentage --- how full planes are on average. One of the seven raw clustering inputs; collapses from $\sim$85\% to $\sim$59\% in 2020, driving the COVID outlier cluster. &
\href{https://en.wikipedia.org/wiki/Load_factor_(aeronautics)}{Wikipedia} \\

PR & Participation Ratio &
A statistic (Eq.~1 in the paper) computed as $(\sum_i \lambda_i)^2 / \sum_i \lambda_i^2$, where $\lambda_i$ are the covariance-matrix eigenvalues. It estimates how many dimensions are ``effectively'' independent versus redundant; borrowed from physics (localization of quantum states) and used in ML as a collinearity/effective-dimensionality diagnostic. Here PR $\approx$ 2.46 of 7 nominal variables. &
\href{https://en.wikipedia.org/wiki/Anderson_localization#Inverse_participation_ratio}{Wikipedia} \\

Sil & Silhouette Coefficient &
For each point, $(b - a)/\max(a,b)$, where $a$ is mean intra-cluster distance and $b$ is mean distance to the nearest other cluster; averaged over all points for the score. Ranges $-1$ to $1$, higher is better. Introduced by Rousseeuw (1987), cited as ref.~[6]. &
\href{https://en.wikipedia.org/wiki/Silhouette_(clustering)}{Wikipedia} \\

DB & Davies--Bouldin Index &
For each cluster, the maximum ratio of (within-cluster scatter of cluster $i$ + within-cluster scatter of cluster $j$) to (distance between centroids $i$ and $j$) over all other clusters $j$, then averaged. Lower values indicate tighter, better-separated clusters. Introduced by Davies \& Bouldin (1979), cited as ref.~[1]. &
\href{https://en.wikipedia.org/wiki/Davies%E2%80%93Bouldin_index}{Wikipedia} \\

ARI & Adjusted Rand Index &
A chance-corrected version of the Rand Index, which counts pairs of points agreeing (both together or both separate) across two partitions. ARI adjusts for chance agreement, so 1.0 = identical partitions and $\approx$0 = random agreement. Introduced by Hubert \& Arabie (1985), cited as ref.~[3]. &
\href{https://en.wikipedia.org/wiki/Rand_index#Adjusted_Rand_index}{Wikipedia} \\

kpca / kPCA & Kernel PCA &
A nonlinear generalization of PCA: instead of eigendecomposing the covariance matrix directly, it eigendecomposes a kernel (Gram) matrix of pairwise similarities computed via a kernel function (linear, RBF, polynomial, etc.), implicitly projecting into a higher-dimensional feature space where linear PCA is then applied. Introduced by Sch\"olkopf, Smola \& M\"uller (1998), cited as ref.~[8]. &
\href{https://en.wikipedia.org/wiki/Kernel_principal_component_analysis}{Wikipedia} \\

RBF & Radial Basis Function (kernel) &
A kernel of the form $K(x,y) = \exp(-\gamma\lVert x-y \rVert^2)$, producing similarity that decays smoothly and symmetrically with Euclidean distance. Widely used as a default nonlinear kernel in kernel PCA, SVMs, and kernel ridge regression because it can approximate arbitrarily smooth functions. &
\href{https://en.wikipedia.org/wiki/Radial_basis_function_kernel}{Wikipedia} \\

$\gamma$ (gamma) & Kernel bandwidth &
The scale parameter in the RBF kernel controlling how quickly similarity decays with distance; small $\gamma$ = smoother/wider similarity, large $\gamma$ = sharper/narrower. Set here via the median pairwise-distance heuristic (a common data-driven default avoiding cross-validation). &
\href{https://scikit-learn.org/stable/modules/generated/sklearn.gaussian_process.kernels.RBF.html}{scikit-learn} \\

5-NN & 5-Nearest Neighbors &
A graph-construction rule where each point is connected only to its 5 closest points by distance, producing a sparse local-connectivity graph rather than a fully dense similarity matrix. Used as the basis for the graph-diffusion kernel, where similarity propagates only along these local edges. &
\href{https://en.wikipedia.org/wiki/K-nearest_neighbors_algorithm}{Wikipedia} \\

FDA & Fisher Discriminant Analysis &
A supervised technique that finds the linear projection maximizing the ratio of between-class variance to within-class variance, given known class labels. The paper argues it's unsuitable here since it requires pre-labelled clusters (circular for structure-discovery) and is numerically unstable with only $\sim$5--9 points per cluster. &
\href{https://en.wikipedia.org/wiki/Linear_discriminant_analysis#Fisher's_linear_discriminant}{Wikipedia} \\

LDA & Linear Discriminant Analysis &
A supervised classification and dimensionality-reduction method closely related to FDA, projecting data to maximize class separability. Mentioned as a possible future validation step once labelled data is available, to check whether supervised projections align with the unsupervised PCA findings. &
\href{https://en.wikipedia.org/wiki/Linear_discriminant_analysis}{Wikipedia} \\

LOOCV & Leave-One-Out Cross-Validation &
A cross-validation scheme where, for $n$ observations, the model is trained $n$ times, each time leaving out exactly one point for testing and training on the rest --- an extreme case of $k$-fold CV with $k=n$. Used with a nested grid search in the kernel ridge regression check. &
\href{https://en.wikipedia.org/wiki/Cross-validation_(statistics)#Leave-one-out_cross-validation}{Wikipedia} \\

R$^2$ & Coefficient of Determination &
The proportion of variance in the dependent variable (here, Operating Profit) explained by the model's predictions, ranging up to 1.0 for a perfect fit (can go negative for a poor model). Reported per model in Table 5, both including and excluding the 2020 outlier year. &
\href{https://en.wikipedia.org/wiki/Coefficient_of_determination}{Wikipedia} \\

HMM & Hidden Markov Model &
A statistical model of sequential data assuming an unobserved (``hidden'') state process that evolves via Markov transitions, with each state emitting observable outputs. Proposed as future work to quantify regime persistence and transition probabilities between the paper's profit-cycle clusters. &
\href{https://en.wikipedia.org/wiki/Hidden_Markov_model}{Wikipedia} \\

VAR & Vector Autoregression &
A multivariate time-series model where each variable is regressed on lagged values of itself and every other variable in the system, up to lag order $p$. Proposed as a ``regime-switching VAR,'' where coefficients depend on a latent regime state (linking to the HMM idea). Would need far more than the current $n=26$ yearly observations given the parameter count scales with $n_{\text{vars}}^2 \times$ lag order. &
\href{https://en.wikipedia.org/wiki/Vector_autoregression}{Wikipedia} \\

t-SNE & t-Distributed Stochastic Neighbor Embedding &
Converts high-dimensional pairwise distances into neighbor probabilities via a Gaussian kernel (tuned by a \textbf{perplexity} parameter setting effective neighborhood size), then fits a low-dimensional layout whose neighbor probabilities --- computed with a heavy-tailed Student's t-distribution --- best match, minimizing KL divergence. The heavy tails avoid the ``crowding problem'' of squeezing high-D neighborhoods into 2D. Purely local/qualitative: axes and inter-cluster distances aren't meaningfully interpretable. &
\href{https://en.wikipedia.org/wiki/T-distributed_stochastic_neighbor_embedding}{Wikipedia} \\

UMAP & Uniform Manifold Approximation and Projection &
Builds a weighted $k$-nearest-neighbor graph approximating the data's local manifold topology (with an adjustable \textbf{n\_neighbors} parameter), then optimizes a low-dimensional embedding via cross-entropy loss to reproduce that fuzzy topological structure. Generally faster than t-SNE and often considered better at preserving some global structure, though results remain stochastic and axes uninterpretable. &
\href{https://umap-learn.readthedocs.io/en/latest/how_umap_works.html}{UMAP docs} \\

\bottomrule
\end{longtable}

\end{document}
